% Studies-in-Algebra-and-Group-Theory - Lecture 12: Generators, Presentations, and Cayley Graphs
% Author: S. D. V. Stephens
% Date: 2026-03-10
% Compile with: pdflatex -shell-escape

\documentclass{report}
\input{~/university/preamble.tex}
% preamble-darkmode.tex
% Dark mode extension for lecture notes
% Usage: % preamble-darkmode.tex
% Dark mode extension for lecture notes
% Usage: % preamble-darkmode.tex
% Dark mode extension for lecture notes
% Usage: \input{~/university/preamble-darkmode.tex} AFTER preamble.tex
%
% This provides:
% - Dark background with light text
% - Material Design color palette
% - Ocean blue gradient colors (p1-p9)
% - Enhanced TikZ/PGFPlots for dark backgrounds
% - qq tcolorbox environment
% - tkz-euclide for geometric constructions

% ===========================================
% DARK MODE PACKAGE
% ===========================================
\usepackage{darkmode}
\enabledarkmode

% ===========================================
% ADDITIONAL PACKAGES
% ===========================================
\usepackage{changepage}
\usepackage{ulem}
\usepackage{colortbl}
\usepackage{tkz-euclide}
\usepackage{capt-of}

% ===========================================
% EXTENDED TIKZ LIBRARIES
% ===========================================
\usetikzlibrary{patterns,3d,backgrounds}
\usepgfplotslibrary{groupplots}

% Upgrade pgfplots compatibility
\pgfplotsset{compat=1.18}

% ===========================================
% OCEAN BLUE GRADIENT (p1-p9)
% ===========================================
\definecolor{p1}{HTML}{caf0f8}
\definecolor{p2}{HTML}{ade8f4}
\definecolor{p3}{HTML}{90e0ef}
\definecolor{p4}{HTML}{48cae4}
\definecolor{p5}{HTML}{00b4d8}
\definecolor{p6}{HTML}{0096c7}
\definecolor{p7}{HTML}{0077b6}
\definecolor{p8}{HTML}{023e8a}
\definecolor{p9}{HTML}{03045e}

% Dark background color
\definecolor{pag}{HTML}{293133}

% ===========================================
% MATERIAL DESIGN COLORS
% ===========================================
% Note: myred, myyellow already defined in main preamble
% These are additional/alternate versions
\definecolor{matred}{HTML}{f44336}
\definecolor{matpink}{HTML}{e81e63}
\definecolor{matpurple}{HTML}{9c27b0}
\definecolor{matdeeppurple}{HTML}{673ab7}
\definecolor{matindigo}{HTML}{3f51b5}
\definecolor{matblue}{HTML}{2196f3}
\definecolor{matlightblue}{HTML}{03a9f4}
\definecolor{matcyan}{HTML}{00bcd4}
\definecolor{matteal}{HTML}{009688}
\definecolor{matgreen}{HTML}{4caf50}
\definecolor{matlightgreen}{HTML}{8bc34a}
\definecolor{matlime}{HTML}{cddc39}
\definecolor{matyellow}{HTML}{ffeb3b}
\definecolor{matamber}{HTML}{ffc107}
\definecolor{matorange}{HTML}{ff9800}
\definecolor{matdeeporange}{HTML}{ff5722}
\definecolor{matbrown}{HTML}{795548}
\definecolor{matgray}{HTML}{9e9e9e}
\definecolor{matbluegray}{HTML}{607d8b}

% ===========================================
% COLORLET FOR TIKZ
% ===========================================
\colorlet{xcol}{blue!60!black}

% ===========================================
% DARK MODE TCOLORBOX
% ===========================================
\newtcolorbox{qq}[2][]{%
    boxrule=0.75pt,
    sharp corners,
    colframe=white,
    colback=pag,
    coltext=white,
    #1,
}

% Dark definition box
\newtcolorbox{darkdfn}[2][]{%
    enhanced,
    boxrule=1pt,
    sharp corners,
    colframe=p5,
    colback=pag,
    coltext=white,
    coltitle=white,
    fonttitle=\bfseries,
    title=#2,
    #1,
}

% Dark theorem box
\newtcolorbox{darkthm}[2][]{%
    enhanced,
    boxrule=1pt,
    sharp corners,
    colframe=matpurple,
    colback=pag,
    coltext=white,
    coltitle=white,
    fonttitle=\bfseries,
    title=#2,
    #1,
}

% ===========================================
% TIKZ 3D FIX
% ===========================================
% Small fix for canvas is xy plane at z
% https://tex.stackexchange.com/a/48776/121799
\makeatletter
\tikzoption{canvas is xy plane at z}[]{%
    \def\tikz@plane@origin{\pgfpointxyz{0}{0}{#1}}%
    \def\tikz@plane@x{\pgfpointxyz{1}{0}{#1}}%
    \def\tikz@plane@y{\pgfpointxyz{0}{1}{#1}}%
    \tikz@canvas@is@plane}
\makeatother

% ===========================================
% CONDITIONAL LINE TIKZSET
% ===========================================
\tikzset{
    conditional line/.style 2 args={
        /utils/exec={\pgfmathsetmacro{\xcoordA}{#1}},
        /utils/exec={\pgfmathsetmacro{\xcoordB}{#2}},
        insert path={
            \ifdim\xcoordA pt>\xcoordB pt
            (\xcoordA,0) -- (\xcoordB,0)
            \fi
        },
    },
}

% ===========================================
% PGFPLOTS DARK MODE DEFAULTS
% ===========================================
\pgfplotsset{
    dark axis/.style={
        axis line style={white},
        tick style={white},
        label style={white},
        grid style={pag!70},
        every axis label/.append style={white},
        every tick label/.append style={white},
    },
}

% ===========================================
% TIKZ EXTERNALIZATION (for fast compilation)
% ===========================================
% This creates .md5 cache files for figures
% Requires: pdflatex -shell-escape
%
% To enable, add AFTER loading this file:
%   \enabletikzexternalize
%
\newcommand{\enabletikzexternalize}{%
  \usetikzlibrary{external}%
  \tikzexternalize[prefix=figures/]%
}

% ===========================================
% CONVENIENT SHORTCUTS FOR DARK PLOTS
% ===========================================
\newcommand{\darkplot}[1]{%
    \begin{tikzpicture}
    \begin{axis}[
        dark axis,
        grid=both,
        major grid style={line width=.2pt,draw=pag!70},
        #1
    ]
}


% ===========================================
% QUICK GEOMETRY MACROS (tkz-euclide)
% ===========================================
% Draw a point: \point{name}{x}{y}
\newcommand{\gpoint}[3]{\tkzDefPoint(#2,#3){#1}}

% Draw line through points: \gline{A}{B}
\newcommand{\gline}[2]{\tkzDrawLine[color=p4](#1,#2)}

% Draw circle: \gcircle{center}{radius}
\newcommand{\gcircle}[2]{\tkzDrawCircle[color=p5](#1,#2)}

% Mark angle: \gangle{A}{B}{C}
\newcommand{\gangle}[3]{\tkzMarkAngle[color=p3,size=0.5](#1,#2,#3)}

% ===========================================
% USAGE EXAMPLE
% ===========================================
% In your document:
%
% \begin{tikzpicture}
%   \begin{axis}[dark axis, ...]
%     \addplot[p4, thick] {sin(deg(x))};
%   \end{axis}
% \end{tikzpicture}
%
% Or with qq box:
% \begin{qq}{My Title}
%   Content here in dark box
% \end{qq}
 AFTER preamble.tex
%
% This provides:
% - Dark background with light text
% - Material Design color palette
% - Ocean blue gradient colors (p1-p9)
% - Enhanced TikZ/PGFPlots for dark backgrounds
% - qq tcolorbox environment
% - tkz-euclide for geometric constructions

% ===========================================
% DARK MODE PACKAGE
% ===========================================
\usepackage{darkmode}
\enabledarkmode

% ===========================================
% ADDITIONAL PACKAGES
% ===========================================
\usepackage{changepage}
\usepackage{ulem}
\usepackage{colortbl}
\usepackage{tkz-euclide}
\usepackage{capt-of}

% ===========================================
% EXTENDED TIKZ LIBRARIES
% ===========================================
\usetikzlibrary{patterns,3d,backgrounds}
\usepgfplotslibrary{groupplots}

% Upgrade pgfplots compatibility
\pgfplotsset{compat=1.18}

% ===========================================
% OCEAN BLUE GRADIENT (p1-p9)
% ===========================================
\definecolor{p1}{HTML}{caf0f8}
\definecolor{p2}{HTML}{ade8f4}
\definecolor{p3}{HTML}{90e0ef}
\definecolor{p4}{HTML}{48cae4}
\definecolor{p5}{HTML}{00b4d8}
\definecolor{p6}{HTML}{0096c7}
\definecolor{p7}{HTML}{0077b6}
\definecolor{p8}{HTML}{023e8a}
\definecolor{p9}{HTML}{03045e}

% Dark background color
\definecolor{pag}{HTML}{293133}

% ===========================================
% MATERIAL DESIGN COLORS
% ===========================================
% Note: myred, myyellow already defined in main preamble
% These are additional/alternate versions
\definecolor{matred}{HTML}{f44336}
\definecolor{matpink}{HTML}{e81e63}
\definecolor{matpurple}{HTML}{9c27b0}
\definecolor{matdeeppurple}{HTML}{673ab7}
\definecolor{matindigo}{HTML}{3f51b5}
\definecolor{matblue}{HTML}{2196f3}
\definecolor{matlightblue}{HTML}{03a9f4}
\definecolor{matcyan}{HTML}{00bcd4}
\definecolor{matteal}{HTML}{009688}
\definecolor{matgreen}{HTML}{4caf50}
\definecolor{matlightgreen}{HTML}{8bc34a}
\definecolor{matlime}{HTML}{cddc39}
\definecolor{matyellow}{HTML}{ffeb3b}
\definecolor{matamber}{HTML}{ffc107}
\definecolor{matorange}{HTML}{ff9800}
\definecolor{matdeeporange}{HTML}{ff5722}
\definecolor{matbrown}{HTML}{795548}
\definecolor{matgray}{HTML}{9e9e9e}
\definecolor{matbluegray}{HTML}{607d8b}

% ===========================================
% COLORLET FOR TIKZ
% ===========================================
\colorlet{xcol}{blue!60!black}

% ===========================================
% DARK MODE TCOLORBOX
% ===========================================
\newtcolorbox{qq}[2][]{%
    boxrule=0.75pt,
    sharp corners,
    colframe=white,
    colback=pag,
    coltext=white,
    #1,
}

% Dark definition box
\newtcolorbox{darkdfn}[2][]{%
    enhanced,
    boxrule=1pt,
    sharp corners,
    colframe=p5,
    colback=pag,
    coltext=white,
    coltitle=white,
    fonttitle=\bfseries,
    title=#2,
    #1,
}

% Dark theorem box
\newtcolorbox{darkthm}[2][]{%
    enhanced,
    boxrule=1pt,
    sharp corners,
    colframe=matpurple,
    colback=pag,
    coltext=white,
    coltitle=white,
    fonttitle=\bfseries,
    title=#2,
    #1,
}

% ===========================================
% TIKZ 3D FIX
% ===========================================
% Small fix for canvas is xy plane at z
% https://tex.stackexchange.com/a/48776/121799
\makeatletter
\tikzoption{canvas is xy plane at z}[]{%
    \def\tikz@plane@origin{\pgfpointxyz{0}{0}{#1}}%
    \def\tikz@plane@x{\pgfpointxyz{1}{0}{#1}}%
    \def\tikz@plane@y{\pgfpointxyz{0}{1}{#1}}%
    \tikz@canvas@is@plane}
\makeatother

% ===========================================
% CONDITIONAL LINE TIKZSET
% ===========================================
\tikzset{
    conditional line/.style 2 args={
        /utils/exec={\pgfmathsetmacro{\xcoordA}{#1}},
        /utils/exec={\pgfmathsetmacro{\xcoordB}{#2}},
        insert path={
            \ifdim\xcoordA pt>\xcoordB pt
            (\xcoordA,0) -- (\xcoordB,0)
            \fi
        },
    },
}

% ===========================================
% PGFPLOTS DARK MODE DEFAULTS
% ===========================================
\pgfplotsset{
    dark axis/.style={
        axis line style={white},
        tick style={white},
        label style={white},
        grid style={pag!70},
        every axis label/.append style={white},
        every tick label/.append style={white},
    },
}

% ===========================================
% TIKZ EXTERNALIZATION (for fast compilation)
% ===========================================
% This creates .md5 cache files for figures
% Requires: pdflatex -shell-escape
%
% To enable, add AFTER loading this file:
%   \enabletikzexternalize
%
\newcommand{\enabletikzexternalize}{%
  \usetikzlibrary{external}%
  \tikzexternalize[prefix=figures/]%
}

% ===========================================
% CONVENIENT SHORTCUTS FOR DARK PLOTS
% ===========================================
\newcommand{\darkplot}[1]{%
    \begin{tikzpicture}
    \begin{axis}[
        dark axis,
        grid=both,
        major grid style={line width=.2pt,draw=pag!70},
        #1
    ]
}


% ===========================================
% QUICK GEOMETRY MACROS (tkz-euclide)
% ===========================================
% Draw a point: \point{name}{x}{y}
\newcommand{\gpoint}[3]{\tkzDefPoint(#2,#3){#1}}

% Draw line through points: \gline{A}{B}
\newcommand{\gline}[2]{\tkzDrawLine[color=p4](#1,#2)}

% Draw circle: \gcircle{center}{radius}
\newcommand{\gcircle}[2]{\tkzDrawCircle[color=p5](#1,#2)}

% Mark angle: \gangle{A}{B}{C}
\newcommand{\gangle}[3]{\tkzMarkAngle[color=p3,size=0.5](#1,#2,#3)}

% ===========================================
% USAGE EXAMPLE
% ===========================================
% In your document:
%
% \begin{tikzpicture}
%   \begin{axis}[dark axis, ...]
%     \addplot[p4, thick] {sin(deg(x))};
%   \end{axis}
% \end{tikzpicture}
%
% Or with qq box:
% \begin{qq}{My Title}
%   Content here in dark box
% \end{qq}
 AFTER preamble.tex
%
% This provides:
% - Dark background with light text
% - Material Design color palette
% - Ocean blue gradient colors (p1-p9)
% - Enhanced TikZ/PGFPlots for dark backgrounds
% - qq tcolorbox environment
% - tkz-euclide for geometric constructions

% ===========================================
% DARK MODE PACKAGE
% ===========================================
\usepackage{darkmode}
\enabledarkmode

% ===========================================
% ADDITIONAL PACKAGES
% ===========================================
\usepackage{changepage}
\usepackage{ulem}
\usepackage{colortbl}
\usepackage{tkz-euclide}
\usepackage{capt-of}

% ===========================================
% EXTENDED TIKZ LIBRARIES
% ===========================================
\usetikzlibrary{patterns,3d,backgrounds}
\usepgfplotslibrary{groupplots}

% Upgrade pgfplots compatibility
\pgfplotsset{compat=1.18}

% ===========================================
% OCEAN BLUE GRADIENT (p1-p9)
% ===========================================
\definecolor{p1}{HTML}{caf0f8}
\definecolor{p2}{HTML}{ade8f4}
\definecolor{p3}{HTML}{90e0ef}
\definecolor{p4}{HTML}{48cae4}
\definecolor{p5}{HTML}{00b4d8}
\definecolor{p6}{HTML}{0096c7}
\definecolor{p7}{HTML}{0077b6}
\definecolor{p8}{HTML}{023e8a}
\definecolor{p9}{HTML}{03045e}

% Dark background color
\definecolor{pag}{HTML}{293133}

% ===========================================
% MATERIAL DESIGN COLORS
% ===========================================
% Note: myred, myyellow already defined in main preamble
% These are additional/alternate versions
\definecolor{matred}{HTML}{f44336}
\definecolor{matpink}{HTML}{e81e63}
\definecolor{matpurple}{HTML}{9c27b0}
\definecolor{matdeeppurple}{HTML}{673ab7}
\definecolor{matindigo}{HTML}{3f51b5}
\definecolor{matblue}{HTML}{2196f3}
\definecolor{matlightblue}{HTML}{03a9f4}
\definecolor{matcyan}{HTML}{00bcd4}
\definecolor{matteal}{HTML}{009688}
\definecolor{matgreen}{HTML}{4caf50}
\definecolor{matlightgreen}{HTML}{8bc34a}
\definecolor{matlime}{HTML}{cddc39}
\definecolor{matyellow}{HTML}{ffeb3b}
\definecolor{matamber}{HTML}{ffc107}
\definecolor{matorange}{HTML}{ff9800}
\definecolor{matdeeporange}{HTML}{ff5722}
\definecolor{matbrown}{HTML}{795548}
\definecolor{matgray}{HTML}{9e9e9e}
\definecolor{matbluegray}{HTML}{607d8b}

% ===========================================
% COLORLET FOR TIKZ
% ===========================================
\colorlet{xcol}{blue!60!black}

% ===========================================
% DARK MODE TCOLORBOX
% ===========================================
\newtcolorbox{qq}[2][]{%
    boxrule=0.75pt,
    sharp corners,
    colframe=white,
    colback=pag,
    coltext=white,
    #1,
}

% Dark definition box
\newtcolorbox{darkdfn}[2][]{%
    enhanced,
    boxrule=1pt,
    sharp corners,
    colframe=p5,
    colback=pag,
    coltext=white,
    coltitle=white,
    fonttitle=\bfseries,
    title=#2,
    #1,
}

% Dark theorem box
\newtcolorbox{darkthm}[2][]{%
    enhanced,
    boxrule=1pt,
    sharp corners,
    colframe=matpurple,
    colback=pag,
    coltext=white,
    coltitle=white,
    fonttitle=\bfseries,
    title=#2,
    #1,
}

% ===========================================
% TIKZ 3D FIX
% ===========================================
% Small fix for canvas is xy plane at z
% https://tex.stackexchange.com/a/48776/121799
\makeatletter
\tikzoption{canvas is xy plane at z}[]{%
    \def\tikz@plane@origin{\pgfpointxyz{0}{0}{#1}}%
    \def\tikz@plane@x{\pgfpointxyz{1}{0}{#1}}%
    \def\tikz@plane@y{\pgfpointxyz{0}{1}{#1}}%
    \tikz@canvas@is@plane}
\makeatother

% ===========================================
% CONDITIONAL LINE TIKZSET
% ===========================================
\tikzset{
    conditional line/.style 2 args={
        /utils/exec={\pgfmathsetmacro{\xcoordA}{#1}},
        /utils/exec={\pgfmathsetmacro{\xcoordB}{#2}},
        insert path={
            \ifdim\xcoordA pt>\xcoordB pt
            (\xcoordA,0) -- (\xcoordB,0)
            \fi
        },
    },
}

% ===========================================
% PGFPLOTS DARK MODE DEFAULTS
% ===========================================
\pgfplotsset{
    dark axis/.style={
        axis line style={white},
        tick style={white},
        label style={white},
        grid style={pag!70},
        every axis label/.append style={white},
        every tick label/.append style={white},
    },
}

% ===========================================
% TIKZ EXTERNALIZATION (for fast compilation)
% ===========================================
% This creates .md5 cache files for figures
% Requires: pdflatex -shell-escape
%
% To enable, add AFTER loading this file:
%   \enabletikzexternalize
%
\newcommand{\enabletikzexternalize}{%
  \usetikzlibrary{external}%
  \tikzexternalize[prefix=figures/]%
}

% ===========================================
% CONVENIENT SHORTCUTS FOR DARK PLOTS
% ===========================================
\newcommand{\darkplot}[1]{%
    \begin{tikzpicture}
    \begin{axis}[
        dark axis,
        grid=both,
        major grid style={line width=.2pt,draw=pag!70},
        #1
    ]
}


% ===========================================
% QUICK GEOMETRY MACROS (tkz-euclide)
% ===========================================
% Draw a point: \point{name}{x}{y}
\newcommand{\gpoint}[3]{\tkzDefPoint(#2,#3){#1}}

% Draw line through points: \gline{A}{B}
\newcommand{\gline}[2]{\tkzDrawLine[color=p4](#1,#2)}

% Draw circle: \gcircle{center}{radius}
\newcommand{\gcircle}[2]{\tkzDrawCircle[color=p5](#1,#2)}

% Mark angle: \gangle{A}{B}{C}
\newcommand{\gangle}[3]{\tkzMarkAngle[color=p3,size=0.5](#1,#2,#3)}

% ===========================================
% USAGE EXAMPLE
% ===========================================
% In your document:
%
% \begin{tikzpicture}
%   \begin{axis}[dark axis, ...]
%     \addplot[p4, thick] {sin(deg(x))};
%   \end{axis}
% \end{tikzpicture}
%
% Or with qq box:
% \begin{qq}{My Title}
%   Content here in dark box
% \end{qq}


\course{Studies-in-Algebra-and-Group-Theory}

\begin{document}

\lecture{12}{Generators, Presentations, and Cayley Graphs}

We now study how to describe groups in terms of generators and relations, and how to visualize their structure via Cayley graphs. This connects algebra to combinatorics and geometry, and leads naturally to the braid group and \(\mathrm{SL}_2(\ZZ)\).


\section{Free Groups and Presentations}

\subsection{The Free Group}

Recall from \textbf{Lecture 3}: the free group \(F_n\) on \(n\) generators \(a_1, \ldots, a_n\) consists of all reduced words \(a_{i_1}^{m_1} a_{i_2}^{m_2} \cdots a_{i_k}^{m_k}\) (where \(k \geq 0\), \(i_j \in \{1, \ldots, n\}\), \(i_j \neq i_{j+1}\), and \(m_j \in \ZZ \setminus \{0\}\)), with multiplication by concatenation followed by reduction (canceling adjacent inverse pairs \(a_i a_i^{-1}\) and combining adjacent powers of the same generator).

\(F_n\) is the ``largest'' group with \(n\) generators: every other group generated by \(n\) elements is a quotient of \(F_n\). If \(G\) is generated by \(g_1, \ldots, g_n \in G\), define a surjective homomorphism \(\varphi : F_n \twoheadrightarrow G\) by \(a_i \mapsto g_i\).

\dfn{Finitely presented group}{A group \(G\) is finitely presented if the kernel of \(\varphi\) is the smallest normal subgroup of \(F_n\) containing some finite set \(\{r_1, \ldots, r_k\} \subset F_n\), i.e., \(\Ker(\varphi) = \langle\!\langle r_1, \ldots, r_k \rangle\!\rangle\) (the normal closure: the subgroup generated by all conjugates \(x r_i x^{-1}\)). We write
\[G \cong \langle a_1, \ldots, a_n \mid r_1, \ldots, r_k \rangle\]
where the \(a_i\) are generators and the \(r_i\) are relations.
}

\ex{Presentations of familiar groups}{
\begin{itemize}
    \item \(\ZZ^n \cong \langle a_1, \ldots, a_n \mid a_i a_j a_i^{-1} a_j^{-1} \; \forall \, i, j \rangle\) (all generators commute).
    \item \(S_3 \cong \langle s_1, s_2 \mid s_1^2, \, s_2^2, \, (s_1 s_2)^3 \rangle\) where \(s_1 = (12)\), \(s_2 = (23)\).
    \item \(S_n \cong \langle s_1, \ldots, s_{n-1} \mid s_i^2 = 1, \; s_i s_j = s_j s_i \; (|i-j| \geq 2), \; s_i s_{i+1} s_i = s_{i+1} s_i s_{i+1} \rangle\) where \(s_i = (i \;\; i\!+\!1)\).
\end{itemize}
}

\nt{Given generators \(g_1, \ldots, g_n \in G\), finding relations \(r_1, \ldots, r_k\) such that these are a complete set (i.e., they generate the full kernel) is nontrivial. If the relations hold in \(G\), then \(\varphi\) induces a surjection \(\langle a_i \mid r_j \rangle \twoheadrightarrow G\). This is an isomorphism iff the relations generate the entire kernel of \(\varphi\).}


\subsection{The Word Problem}

\dfn{Word problem}{Given a presentation \(G = \langle a_1, \ldots, a_n \mid r_1, \ldots, r_k \rangle\), the word problem asks: given a word \(w\) in the generators, does \(w\) represent the identity in \(G\)?}

\nt{For groups where we can ``calculate'' (e.g., matrix groups, permutation groups), the word problem is trivially solvable: just evaluate the word. But for abstractly presented groups, this can be undecidable. Novikov (1955) and Boone (1959) independently showed there exist finitely presented groups with unsolvable word problem. Two useful tools for groups where the word problem is solvable: Cayley graphs and normal forms.}


\section{Cayley Graphs}

\dfn{Cayley graph}{Given a group \(G\) with generators \(g_1, \ldots, g_n\), the Cayley graph of \(G\) has:
\begin{itemize}
    \item Vertices: the elements of \(G\).
    \item Edges: for each \(s \in G\) and generator \(g_i\), a directed edge from \(s\) to \(sg_i\), labeled \(g_i\).
\end{itemize}
}

\nt{\(G\) acts on its Cayley graph by left multiplication: \(g\) sends vertex \(s\) to \(gs\) and edge \((s \to sg_i)\) to \((gs \to gsg_i)\). This action is transitive on vertices and preserves the graph structure (it is a graph automorphism). The graph is highly symmetric.}

\ex{Cayley graph of \(\ZZ\)}{With generator \(1\), the Cayley graph is the integer number line:
\begin{center}
\begin{tikzpicture}[scale=1.2]
\foreach \x in {-3,...,3} {
    \node[circle, fill=black, inner sep=1.5pt, label=below:{\small \(\x\)}] (n\x) at (\x, 0) {};
}
\foreach \x in {-3,...,2} {
    \pgfmathtruncatemacro{\y}{\x+1}
    \draw[->, >=stealth, thick] (n\x) -- (n\y);
}
\draw[dotted, thick] (-3.5, 0) -- (-3, 0);
\draw[dotted, thick] (3, 0) -- (3.5, 0);
\end{tikzpicture}
\end{center}
}

\ex{Cayley graph of \(S_3\)}{With generators \(s_1 = (12)\) and \(s_2 = (23)\): the 6 vertices form a hexagonal graph. The relation \((s_1 s_2)^3 = e\) corresponds to the outer hexagonal cycle closing up.
\begin{center}
\begin{tikzpicture}[scale=1.6, every node/.style={circle, draw, fill=white, inner sep=2pt, minimum size=18pt, font=\footnotesize}]
    \node (e)   at (90:1.8)  {\(e\)};
    \node (s1)  at (150:1.8) {\(s_1\)};
    \node (s2)  at (30:1.8)  {\(s_2\)};
    \node (s1s2) at (210:1.8) {\(s_1 s_2\)};
    \node (s2s1) at (330:1.8) {\(s_2 s_1\)};
    \node (s1s2s1) at (270:1.8) {\(s_1 s_2 s_1\)};
    \draw[->, >=stealth, thick, blue] (e) -- node[draw=none, fill=none, above left, font=\tiny] {\(s_1\)} (s1);
    \draw[->, >=stealth, thick, red] (e) -- node[draw=none, fill=none, above right, font=\tiny] {\(s_2\)} (s2);
    \draw[->, >=stealth, thick, red] (s1) -- node[draw=none, fill=none, left, font=\tiny] {\(s_2\)} (s1s2);
    \draw[->, >=stealth, thick, blue] (s2) -- node[draw=none, fill=none, right, font=\tiny] {\(s_1\)} (s2s1);
    \draw[->, >=stealth, thick, blue] (s1s2) -- node[draw=none, fill=none, below left, font=\tiny] {\(s_1\)} (s1s2s1);
    \draw[->, >=stealth, thick, red] (s2s1) -- node[draw=none, fill=none, below right, font=\tiny] {\(s_2\)} (s1s2s1);
\end{tikzpicture}
\end{center}
Since any word in \(s_1, s_2\) with \(s_1^2 = s_2^2 = e\) reduces to an alternating sequence \((\cdots s_1 s_2 s_1 \cdots)\), we verify \(S_3 \cong \langle s_1, s_2 \mid s_1^2, s_2^2, (s_1 s_2)^3 \rangle\).
}

\ex{Cayley graph of \(S_4\)}{With generators \(s_i = (i \;\; i\!+\!1)\) for \(i = 1, 2, 3\): the 24 vertices form the edge graph of a truncated octahedron (a permutohedron). The faces correspond to the relations: hexagonal faces to \(s_i s_{i+1} s_i = s_{i+1} s_i s_{i+1}\) and square faces to \(s_1 s_3 = s_3 s_1\).
}

\dfn{Word length}{The word length of \(g \in G\) (with respect to generators \(g_1, \ldots, g_n\)) is the shortest distance from \(e\) to \(g\) in the Cayley graph, i.e., the minimal number of generators or their inverses needed to express \(g\).
}

\nt{For \(S_n\) with generators \(s_i = (i \;\; i\!+\!1)\), the word length of a permutation \(\sigma\) equals the number of inversions \(|\{(i, j) : i < j, \, \sigma(i) > \sigma(j)\}|\). The longest element is \(\sigma = \binom{1 \; 2 \; \cdots \; n}{n \; n\!-\!1 \; \cdots \; 1}\) with word length \(\binom{n}{2}\).}

\subsection{Growth Rate}

For infinite groups, the Cayley graph encodes asymptotic information about the group.

\dfn{Growth rate}{Given generators \(g_1, \ldots, g_n\) of an infinite group \(G\), the growth function counts \(B(N) = |\{g \in G : \text{word length of } g \leq N\}|\). The group has polynomial growth if \(B(N) \sim N^d\) for some \(d\), and exponential growth if \(B(N) \sim c^N\) for some \(c > 1\).
}

\nt{The growth type (polynomial vs exponential) is independent of the choice of generators: changing generators changes word lengths by at most a bounded factor (each new generator is a bounded-length word in the old generators and vice versa). Finitely generated abelian groups have polynomial growth. Free groups have exponential growth. Gromov's theorem (1981) characterizes polynomial growth: a finitely generated group has polynomial growth iff it is virtually nilpotent (has a nilpotent subgroup of finite index).}


\section{The Braid Group}

\dfn{Braid group}{The braid group \(B_n\) on \(n\) strands has presentation
\[B_n = \langle s_1, \ldots, s_{n-1} \mid s_i s_j = s_j s_i \; (|i-j| \geq 2), \; s_i s_{i+1} s_i = s_{i+1} s_i s_{i+1} \rangle.\]
This is the same as the presentation of \(S_n\) but without the relations \(s_i^2 = 1\). The generator \(s_i\) represents crossing strand \(i\) over strand \(i+1\).

\begin{center}
\begin{tikzpicture}
\pic[braid/.cd, number of strands=3, line width=1.2pt, gap=0.08] at (0,0) {braid={s_1}};
\node at (1, -2.3) {\(s_1\)};
\pic[braid/.cd, number of strands=3, line width=1.2pt, gap=0.08] at (4,0) {braid={s_1^{-1}}};
\node at (5, -2.3) {\(s_1^{-1}\)};
\pic[braid/.cd, number of strands=3, line width=1.2pt, gap=0.08] at (8,0) {braid={s_2}};
\node at (9, -2.3) {\(s_2\)};
\end{tikzpicture}
\end{center}
}

\nt{There is a natural surjection \(B_n \twoheadrightarrow S_n\) sending each generator \(s_i\) to the transposition \((i \;\; i\!+\!1)\). The kernel is the pure braid group \(P_n\). Unlike \(S_n\), the braid group is infinite: \(s_i^2 \neq e\) in \(B_n\).

Braids are best understood as diagrams: \(n\) strands running from top to bottom, where \(s_i\) crosses strand \(i\) over strand \(i+1\), and \(s_i^{-1}\) crosses strand \(i+1\) over strand \(i\). Composition is vertical stacking. Two braid diagrams represent the same element iff one can be deformed into the other (isotopy). The braid relation \(s_1 s_2 s_1 = s_2 s_1 s_2\) (the Yang--Baxter equation) says these two braids are isotopic:

\begin{center}
\begin{tikzpicture}
\pic[braid/.cd, number of strands=3, line width=1.2pt, gap=0.08] at (0,0) {braid={s_1 s_2 s_1}};
\node at (1, -4.8) {\(s_1 s_2 s_1\)};
\node at (4, -2.5) {\(=\)};
\pic[braid/.cd, number of strands=3, line width=1.2pt, gap=0.08] at (5.5,0) {braid={s_2 s_1 s_2}};
\node at (6.5, -4.8) {\(s_2 s_1 s_2\)};
\end{tikzpicture}
\end{center}
}

\nt{Markov's theorem connects braids to knot theory: every knot or link in \(\RR^3\) can be represented as the closure of a braid, and two braids give isotopic links iff they are related by conjugation in \(B_n\) and stabilization \(\sigma \in B_n \sim \sigma s_n^{\pm 1} \in B_{n+1}\).}

\nt{\(B_n\) is much larger than \(S_n\), but its algorithmic aspects can be approached by the same methods, using permutation braids (braids where any two strands cross at most once, in bijection with \(S_n\)). Let \(\Delta\) be the longest permutation braid. Since its shortest word can start/end with any \(s_i\), and \(s_i^{\pm 1} \Delta\) is still a permutation braid, every element of \(B_n\) can be written as \(g = \Delta^k P_1 \cdots P_r\) where the \(P_j\) are permutation braids satisfying a ``left greedy'' condition. This is the Garside normal form (Garside 1969, Thurston--ElRifai--Morton early 1990s), which solves the word problem in \(B_n\).}


\section{\(\mathrm{SL}_2(\ZZ)\) and \(\mathrm{PSL}_2(\ZZ)\)}

\dfn{\(\mathrm{SL}_2(\ZZ)\)}{\(\mathrm{SL}_2(\ZZ) = \left\{\begin{pmatrix} a & b \\ c & d \end{pmatrix} : a, b, c, d \in \ZZ, \; ad - bc = 1\right\}\) is the group of \(2 \times 2\) integer matrices with determinant 1. Its quotient by the center \(\{\pm I\}\) is \(\mathrm{PSL}_2(\ZZ) = \mathrm{SL}_2(\ZZ) / \{\pm I\}\).
}

\mprop{Generators of \(\mathrm{SL}_2(\ZZ)\)}{\(\mathrm{SL}_2(\ZZ)\) is generated by \(S = \begin{pmatrix} 0 & -1 \\ 1 & 0 \end{pmatrix}\) and \(T = \begin{pmatrix} 1 & 1 \\ 0 & 1 \end{pmatrix}\).
}

\pf{Proof}{Given \(M = \begin{pmatrix} a & b \\ c & d \end{pmatrix} \in \mathrm{SL}_2(\ZZ)\), we express \(M\) in terms of \(S\) and \(T\) by repeatedly reducing the entries.

If \(c = 0\): then \(ad = 1\), so \(a = d = \pm 1\). Thus \(M = \pm \begin{pmatrix} 1 & n \\ 0 & 1 \end{pmatrix} = \pm T^n\), which is \(T^n\) or \(S^2 T^n\) (since \(S^2 = -I\)).

If \(c \neq 0\): apply the Euclidean algorithm to reduce \(\max(|a|, |c|)\).
\begin{itemize}
    \item If \(|a| \geq |c|\): write \(a = nc + r\) with \(|r| < |c|\). Then \(T^{-n} M = \begin{pmatrix} a - nc & b - nd \\ c & d \end{pmatrix} = \begin{pmatrix} r & * \\ c & d \end{pmatrix}\), reducing \(\max(|a|, |c|)\).
    \item If \(|a| < |c|\): apply \(S^{-1} M = \begin{pmatrix} c & d \\ -a & -b \end{pmatrix}\), swapping \(a\) and \(c\) (up to sign), then return to the previous case.
\end{itemize}
After finitely many steps, the product of \(M\) with some word in \(S\) and \(T\) has \(c = 0\), hence is \(\pm T^n\). Inverting gives \(M\) as a word in \(S, T\).
}

\nt{Alternative generators: \(S\) and \(R = ST = \begin{pmatrix} 0 & -1 \\ 1 & 1 \end{pmatrix}\) have finite order: \(S^4 = I\) and \(R^6 = I\) (so their images in \(\mathrm{PSL}_2(\ZZ)\) have orders 2 and 3). Also \(T' = \begin{pmatrix} 1 & 0 \\ -1 & 1 \end{pmatrix} = (TST)^{-1} = STS^{-1}\), showing \(T\) and \(T'\) are conjugate.}

\thm{Presentation of \(\mathrm{PSL}_2(\ZZ)\)}{\(\mathrm{PSL}_2(\ZZ) \cong \langle S, R \mid S^2, R^3 \rangle \cong \ZZ/2 * \ZZ/3\), the free product of \(\ZZ/2\) and \(\ZZ/3\).}

\nt{This can be proved geometrically. \(\mathrm{PSL}_2(\ZZ)\) acts on the upper half-plane \(\mathbb{H} = \{z \in \CC : \operatorname{Im}(z) > 0\}\) by Möbius transformations: \(\begin{pmatrix} a & b \\ c & d \end{pmatrix} \cdot z = \frac{az + b}{cz + d}\). Under this action, \(S\) sends \(z \mapsto -1/z\) and \(T\) sends \(z \mapsto z + 1\).

The region \(\Delta = \{z \in \mathbb{H} : |z| \geq 1, \, |\operatorname{Re}(z)| \leq 1/2\}\) is a fundamental domain: its translates under \(\mathrm{PSL}_2(\ZZ)\) tile \(\mathbb{H}\) without overlap.

\begin{center}
\begin{tikzpicture}[scale=2.5]
    % Axes
    \draw[->] (-2.2, 0) -- (2.2, 0) node[right] {\small\(\operatorname{Re}\)};
    \draw[->] (0, -0.05) -- (0, 2.2) node[above] {\small\(\operatorname{Im}\)};
    % Unit circles (arcs) for various translates
    \foreach \k in {-2,-1,0,1,2} {
        \draw[gray, thin] ({\k-0.5}, 0) arc (180:0:0.5);
    }
    \foreach \k in {-1,0,1} {
        \draw[gray, thin] ({\k-1}, 0) arc (180:0:1);
    }
    % Fundamental domain (shaded)
    \fill[blue!15] (-0.5, {sqrt(3)/2}) -- (-0.5, 2.1) -- (0.5, 2.1) -- (0.5, {sqrt(3)/2}) arc (60:120:1) -- cycle;
    \draw[blue, thick] (-0.5, {sqrt(3)/2}) -- (-0.5, 2.1);
    \draw[blue, thick] (0.5, {sqrt(3)/2}) -- (0.5, 2.1);
    \draw[blue, thick] (-0.5, {sqrt(3)/2}) arc (120:60:1);
    % Labels
    \node[blue] at (0, 1.4) {\small\(\Delta\)};
    \node[below] at (-0.5, 0) {\tiny\(-\frac{1}{2}\)};
    \node[below] at (0.5, 0) {\tiny\(\frac{1}{2}\)};
    \node[below] at (-1, 0) {\tiny\(-1\)};
    \node[below] at (1, 0) {\tiny\(1\)};
    % Key labels for adjacent regions
    \node[gray] at (-1, 0.7) {\tiny\(T^{-1}\Delta\)};
    \node[gray] at (1, 0.7) {\tiny\(T\Delta\)};
    \node[gray] at (0, 0.45) {\tiny\(S\Delta\)};
\end{tikzpicture}
\end{center}

The regions adjacent to \(\Delta\) are \(T^{\pm 1}(\Delta)\) and \(S(\Delta)\), and the tiling structure is the Cayley graph of \(\mathrm{PSL}_2(\ZZ)\) with generators \(S, T\). That all regions can be reached from \(\Delta\) in finitely many steps is equivalent to \(S, T\) generating \(\mathrm{PSL}_2(\ZZ)\).}


\mer{Generators, presentations, and Cayley graphs exercises}{
\begin{enumerate}[(1)]
    \item Verify the presentation \(S_3 \cong \langle s_1, s_2 \mid s_1^2, s_2^2, (s_1 s_2)^3 \rangle\) by checking that the 6 reduced words \(e, s_1, s_2, s_1 s_2, s_2 s_1, s_1 s_2 s_1\) are all distinct and form a group under the given relations.
    
    \item Show that the braid group \(B_2\) is isomorphic to \(\ZZ\). What is \(B_3\)?
    
    \item Compute \(S^2\), \(S^4\), \((ST)^3\), and \((ST)^6\) in \(\mathrm{SL}_2(\ZZ)\). Use this to verify \(S^4 = I\) and \(R^6 = I\) where \(R = ST\).
    
    \item Express \(M = \begin{pmatrix} 3 & 5 \\ 1 & 2 \end{pmatrix}\) as a word in \(S\) and \(T\) using the Euclidean algorithm from the proof.
\end{enumerate}
}

\begin{solution}
\textbf{(1)} The relations \(s_1^2 = s_2^2 = e\) and \((s_1 s_2)^3 = e\) (i.e., \(s_1 s_2 s_1 = s_2 s_1 s_2\)) mean any word reduces to alternating sequences of bounded length. The 6 reduced words are: \(e, s_1, s_2, s_1 s_2, s_2 s_1, s_1 s_2 s_1\). Compute the multiplication table using the relations: e.g., \((s_1 s_2)(s_2 s_1) = s_1 (s_2^2) s_1 = s_1^2 = e\), so \(s_1 s_2\) and \(s_2 s_1\) are inverses. Check \(s_1 s_2 s_1 = s_2 s_1 s_2\) confirms the 6 elements close under multiplication. Since \(|S_3| = 6\) and the surjection \(\langle s_1, s_2 \mid \cdots \rangle \twoheadrightarrow S_3\) has at most 6 elements in the domain, it is an isomorphism.

\textbf{(2)} \(B_2 = \langle s_1 \rangle\) with no relations (since there are no pairs \(|i - j| \geq 2\) and no braid relations). So \(B_2 \cong \ZZ\). For \(B_3 = \langle s_1, s_2 \mid s_1 s_2 s_1 = s_2 s_1 s_2 \rangle\): this is an infinite non-abelian group. Its center is \(\ZZ\), generated by \((s_1 s_2)^3 = (s_1 s_2 s_1)^2\), and \(B_3 / Z(B_3) \cong \mathrm{PSL}_2(\ZZ)\).

\textbf{(3)} \(S^2 = \begin{pmatrix} 0 & -1 \\ 1 & 0 \end{pmatrix}^2 = \begin{pmatrix} -1 & 0 \\ 0 & -1 \end{pmatrix} = -I\). So \(S^4 = I\).

\(R = ST = \begin{pmatrix} 0 & -1 \\ 1 & 1 \end{pmatrix}\). \(R^2 = \begin{pmatrix} -1 & -1 \\ 1 & 0 \end{pmatrix}\). \(R^3 = R \cdot R^2 = \begin{pmatrix} -1 & 0 \\ 0 & -1 \end{pmatrix} = -I\). So \(R^6 = I\).

In \(\mathrm{PSL}_2(\ZZ)\): \(\bar{S}^2 = \bar{R}^3 = \bar{I}\), confirming orders 2 and 3.

\textbf{(4)} Start: \(M = \begin{pmatrix} 3 & 5 \\ 1 & 2 \end{pmatrix}\). Since \(|a| = 3 > |c| = 1\), write \(3 = 3 \cdot 1 + 0\). Apply \(T^{-3} M = \begin{pmatrix} 1 & -3 \\ 0 & 1 \end{pmatrix} \begin{pmatrix} 3 & 5 \\ 1 & 2 \end{pmatrix} = \begin{pmatrix} 0 & -1 \\ 1 & 2 \end{pmatrix}\).

Now \(|a| = 0 < |c| = 1\), apply \(S^{-1}\): \(S^{-1} \begin{pmatrix} 0 & -1 \\ 1 & 2 \end{pmatrix} = \begin{pmatrix} 0 & 1 \\ -1 & 0 \end{pmatrix} \begin{pmatrix} 0 & -1 \\ 1 & 2 \end{pmatrix} = \begin{pmatrix} 1 & 2 \\ 0 & 1 \end{pmatrix} = T^2\).

So \(S^{-1} T^{-3} M = T^2\), giving \(M = T^3 S T^2\).

Check: \(T^3 S T^2 = \begin{pmatrix} 1 & 3 \\ 0 & 1 \end{pmatrix} \begin{pmatrix} 0 & -1 \\ 1 & 0 \end{pmatrix} \begin{pmatrix} 1 & 2 \\ 0 & 1 \end{pmatrix} = \begin{pmatrix} 1 & 3 \\ 0 & 1 \end{pmatrix} \begin{pmatrix} 0 & -1 \\ 1 & 2 \end{pmatrix} = \begin{pmatrix} 3 & 5 \\ 1 & 2 \end{pmatrix}\). Confirmed.
\end{solution}



\pf{Proof of the presentation}{We show \(\mathrm{PSL}_2(\ZZ) \cong \langle S, R \mid S^2, R^3 \rangle\). We have \(S^2 = -I\) and \(R^3 = -I\) in \(\mathrm{SL}_2(\ZZ)\), so \(\bar{S}\) and \(\bar{R}\) have orders 2 and 3 in \(\mathrm{PSL}_2(\ZZ)\). The relations \(S^2 = R^3 = e\) reduce any word in \(S^{\pm 1}, R^{\pm 1}\) to the form \(\cdots S R^{\pm 1} S R^{\pm 1} \cdots\) (alternating, since the word can start and end with either \(S\) or \(R^{\pm 1}\)).

Map \(F_2 = \langle s, r \rangle \to \mathrm{PSL}_2(\ZZ)\) by \(r \mapsto R\), \(s \mapsto S\). This induces a surjection \(\langle s, r \mid s^2, r^3 \rangle \twoheadrightarrow \mathrm{PSL}_2(\ZZ)\). The kernel consists of all words \(\cdots s r^{\pm 1} s r^{\pm 1} \cdots\) that simplify to \(e\) in \(\mathrm{PSL}_2(\ZZ)\), i.e., to \(\pm I\) in \(\mathrm{SL}_2(\ZZ)\).

Assume some word \(w\) in \(S\) and \(R^{\pm 1}\) (alternating) simplifies to \(\pm I\). If \(w\) starts and ends with \(S\), conjugate by \(S\) to get a shorter word \(w' = SwS\) with \(w' = \pm I\) as well. If \(w\) starts with \(R^{\pm 1}\), conjugate by \(R^{\mp 1}\) to get another word that doesn't. Iterating, we eventually get a word \(SR^{\pm 1} \cdots SR^{\pm 1} = \pm I\).

But \(SR = -T = -\begin{pmatrix} 1 & 1 \\ 0 & 1 \end{pmatrix}\) and \(SR^{-1} = \begin{pmatrix} 1 & 0 \\ 1 & 1 \end{pmatrix}\). So a product \((SR^{\varepsilon_1})(SR^{\varepsilon_2}) \cdots (SR^{\varepsilon_k})\) with each \(\varepsilon_i = \pm 1\) equals \((-1)^j\) (where \(j\) counts the factors with \(\varepsilon_i = +1\)) times a product of the nonnegative matrices \(T = \begin{pmatrix} 1 & 1 \\ 0 & 1 \end{pmatrix}\) and \(\begin{pmatrix} 1 & 0 \\ 1 & 1 \end{pmatrix}\). Any nonempty product of these nonnegative matrices has entry sum \(\geq 3 > 2\) (one can check the sum is non-decreasing and the initial matrices already have sum 3), so it can never equal \(\pm I\) (which has entry sum 2). Thus no nontrivial alternating word equals \(\pm I\), and the kernel is trivial.
}

\nt{Alternative generators and presentations: using \(T = \begin{pmatrix} 1 & 1 \\ 0 & 1 \end{pmatrix}\) and \(T' = (TST)^{-1} = STS^{-1} = \begin{pmatrix} 1 & 0 \\ -1 & 1 \end{pmatrix}\) (which are conjugate), we get \(\mathrm{SL}_2(\ZZ) \cong \langle T, T' \mid (TT')^6 = 1, \, TT'T = T'TT' \rangle\). The relation \(TT'T = T'TT'\) is precisely the braid relation.}

\nt{This connects to our discussion of braid groups: the center of \(B_3 = \langle s_1, s_2 \mid s_1 s_2 s_1 = s_2 s_1 s_2 \rangle\) is generated by \(\Delta^2 = (s_1 s_2)^3\) and is isomorphic to \(\ZZ\). The map \(s_1 \mapsto T\), \(s_2 \mapsto T'\) gives \(1 \to \ZZ = \langle \Delta^2 \rangle \hookrightarrow B_3 \twoheadrightarrow \mathrm{PSL}_2(\ZZ) \to 1\), a central extension.}


\end{document}